% ==========================================
% SECCIÓN 0: INTRODUCCIÓN
% ==========================================
% TIP: Esta sección introduce al lector al tema de la práctica, explicando su objetivo, los conceptos clave que se abordarán (tipos de datos, procesos, archivos), y cómo estos conceptos se relacionan entre sí en el contexto de la programación en C. También se puede incluir una breve historia de C y su importancia en la programación.
%
% TIP: Las rutas de imágenes son relativas a main.tex.
%
% ==========================================

\chapter{Introducción}

El cómputo paralelo constituye una estrategia fundamental para mejorar el desempeño de aplicaciones que requieren procesar grandes volúmenes de información en tiempos reducidos. En este contexto, la programación en lenguaje C y el uso de procesos en sistemas tipo Unix permiten implementar soluciones concurrentes con control detallado de recursos, memoria y comunicación entre tareas. Por ello, comprender la relación entre procesos, tipos de datos y manejo de archivos resulta indispensable para desarrollar programas correctos, eficientes y verificables.

La problemática abordada en esta práctica se centra en implementar un programa que distribuya el trabajo entre procesos hijo, de modo que cada uno ejecute una parte específica del cálculo sobre datos provenientes de un archivo CSV. Este enfoque exige no solo crear y sincronizar procesos de forma adecuada, sino también garantizar que la lectura de datos, el procesamiento y la escritura de resultados se realicen de manera consistente. En consecuencia, el reto principal no se limita al resultado final, sino a la correcta coordinación entre componentes concurrentes.

El objetivo general del reporte es documentar el fundamento teórico y la solución implementada para resolver el problema planteado mediante programación paralela basada en procesos. Además, se consideran la descripción de los tipos de datos utilizados, las llamadas al sistema para gestión de procesos, el tratamiento de archivos y la validación de la ejecución del programa. Asimismo, se presentan evidencias de funcionamiento y una reflexión crítica sobre los aprendizajes obtenidos durante el desarrollo de la práctica.

La estructura del documento se organiza en capítulos que cubren progresivamente los elementos necesarios para comprender la solución. Primero se desarrolla el marco teórico sobre tipos de datos, procesos y archivos; posteriormente se expone el diseño y funcionamiento del programa implementado; finalmente se integran las conclusiones individuales y colectivas del equipo.

% ------------------------------------------
% Responsable de la introducción: Bustillos Cruz Jonatan
% ------------------------------------------

% TODO (Jonatan): Escribir una introducción general sobre la práctica, explicando su objetivo, los conceptos clave que se abordarán (tipos de datos, procesos, archivos), y cómo estos conceptos se relacionan entre sí en el contexto de la programación en C.
% TODO (Jonatan): Incluir una breve historia de C y su importancia en la programación